% !TeX root = ../教材题库_全部条目_转写工作稿.tex
% Source: XJ-B1-C04-RV-P152-Y0638
% 第4章 / 章末复习 复习题四 / 教材 p.152 / PDF 第159页
% OCR未校对：公式与文字均可能有误；未完成数学TeX转写；禁止直接发给学生
% 原题及相邻原解截图（核对用）：../题目原图/XJ-B1-C04-RV-P152-Y0638-01.webp
% 本文件只在转写完成后作为题干使用；OCR 条目尚不是可靠题干。
\begin{bankitem}{复习题四 第15题}{第4章 / 章末复习 复习题四 / 教材 p.152 / PDF 第159页}{XJ-B1-C04-RV-P152-Y0638}{OCR未校对：公式与文字均可能有误；未完成数学TeX转写；禁止直接发给学生}
\begin{ocrdraft}
[教材p.152 / PDF第159页]
\par
巧．
\par
借助计算器或计算机，用二分法求函数不到一分十2工一1在区间@，1]
\par
内的零点的近似值（误差不超过0．1）．
\end{ocrdraft}
\end{bankitem}
